\section{Expanded Related Work}\label{sec:exp_related}


\noindent \textbf{Gradient-Based Agentic Optimization Methods.} As previously discussed, prior research has explored several gradient-based optimization methods for agent systems, such as supervised fine-tuning \cite{api_bank} and prompt-tuning techniques \cite{dsp}. However, most of these approaches are designed for single-agent optimization and do not extend naturally to MAS that require multi-step collaborative interactions among multiple agents. To the best of our knowledge, research on gradient-based methods specifically targeting MAS optimization in multi-step collaboration settings remains very limited.

In scenarios involving multiple interacting agents, gradient-based optimization becomes particularly challenging due to the presence of numerous confounding variables and the difficulty of jointly performance optimization with multiple functional and structural components. For instance, while some works have explored the use of reinforcement learning for optimizing agentic systems, these methods typically focus on single-agent settings \cite{shao2024deepseekmath, qian2025toolrl}, which differ fundamentally from the multi-agent optimization problem addressed in our study.
A few very recent works, such as Spiral \cite{liu2025spiral} and MHGPO \cite{chen2025heterogeneous}, have made initial attempts at applying RL to MAS optimization. However, their formulations are generally constrained to single-interaction and shared-policy scenarios, where all agents share the same control policy and perform only one round of communication or problem-solving. This setup is considerably simpler and less flexible than the setting studied in our work, which involves multi-step collaboration among role-specialized agents operating within dynamic interaction structures. Consequently, these approaches are not directly comparable to OMAC, which provides a general and extensible optimization framework for optimizing multi-step, role-specialized collaborations among multiple agents.



\noindent \textbf{Non-Gradient-Based MAS Optimization Methods.} A very recent work, ADAS \cite{hu2024automated}, takes an initial step toward applying prompt evolution techniques combined with search tools to optimize MAS. However, ADAS focuses on agent functionality optimization, leaving the collaboration structure predefined and fixed. This contrasts with OMAC, which introduces a general framework capable of comprehensively optimizing both the functional and structural aspects of MAS. Beyond ADAS, a couple of recent studies have explored non-gradient-based optimization for MAS from other perspectives. For example, AFlow \cite{aflow} leverages supervised signals to guide Monte Carlo Tree Search (MCTS) for discovering effective agentic workflows. 
In contrast, OMAC uses supervised signals directly for contrastive reasoning–based optimization, representing a distinct paradigm that avoids reliance on search-based methods. Similarly, other recent works such as G-Designer \cite{zhang2025g} and MaAS \cite{zhang2025multi} focus on architectural search for identifying effective multi-agent structures, thereby addressing primarily the structural optimization dimension. 
For example, in the MaAS implementation, the authors directly adopt well-optimized and near-optimal agent configurations from prior work like AgentCoder~\cite{huang2023agentcoder} and ReAct~\cite{yao2023react} as fixed functional components, and only optimize the collaboration structure. OMAC, however, begins with simple and minimally engineered agent configurations and is designed to optimize both agent functionality and MAS structure. 
%Because MaAS relies on externally optimized functional agents and does not address functional optimization itself, we believe it is not directly comparable to OMAC in terms of scope or methodology.
In summary, while these works have advanced non-gradient MAS optimization from various angles, they typically address only a single aspect, either functional or structural, or use supervised signals indirectly within evolutionary search algorithms like MCTS. By contrast, OMAC provides a unified framework that jointly optimizes both functionality and collaboration structure, employing direct contrastive reasoning with supervised signals rather than relying on evolutionary or search-based methods.